
\documentclass{article}
\usepackage{colortbl}
\usepackage{makecell}
\usepackage{multirow}
\usepackage{supertabular}

\begin{document}

\newcounter{utterance}

\centering \large Interaction Transcript for game `dond', experiment `semi\_en', episode 6 with reflexion{-}Qwen3.5{-}27B.
\vspace{24pt}

{ \footnotesize  \setcounter{utterance}{1}
\setlength{\tabcolsep}{0pt}
\begin{supertabular}{c@{$\;$}|p{.15\linewidth}@{}p{.15\linewidth}p{.15\linewidth}p{.15\linewidth}p{.15\linewidth}p{.15\linewidth}}
    \# & $\;$A & \multicolumn{4}{c}{Game Master} & $\;\:$B\\
    \hline

    \theutterance \stepcounter{utterance}  
    & & \multicolumn{4}{p{0.6\linewidth}}{
        \cellcolor[rgb]{0.9,0.9,0.9}{
            \makecell[{{p{\linewidth}}}]{
                \texttt{\tiny{[P1$\langle$GM]}}
                \texttt{You are playing a negotiation game in which you have to agree on how to divide a set of items among you and another player.} \\
\\ 
\texttt{Rules:} \\
\texttt{(a) You and the other player are given a set of items. Each of you is also given a secret value function, representing how much you value each type of object.} \\
\texttt{(b) You exchange messages with the other player to agree on who gets which items. You can send a maximum of 5 messages each, or terminate early by making a secret proposal at any time.} \\
\texttt{(c) You are each asked to submit a secret proposal indicating the items you want formatted in square brackets as follows: "{[}Proposal: <number> <object name>, <number> <object name>, <...>{]}"} \\
\texttt{(d) If your proposals are complementary, i.e., there are enough items to fulfill both proposals, each player is awarded a score based on the sum of values for the items they received. Otherwise, both of you get zero points.} \\
\texttt{(e) Your goal is to maximize the score you receive.} \\
\\ 
\texttt{Let us start.} \\
\\ 
\texttt{The set of available items is:} \\
\\ 
\texttt{2 eggs, 1 picture, 2 dice} \\
\\ 
\texttt{Your secret value function is:} \\
\\ 
\texttt{egg: 1, picture: 2, dice: 3} \\
\\ 
\texttt{IMPORTANT: Your messages, unless it is the secret proposal, are directly transmitted to the other player, so do not include any response to the rules or text announcing your message. To make a secret proposal, use the indicated format. Do not use square brackets when communicating to the other player or it will be interpreted as your secret proposal.} \\
            }
        }
    }
    & \\ \\

    \theutterance \stepcounter{utterance}  
    & \multicolumn{4}{p{0.6\linewidth}}{
        \cellcolor[rgb]{0.9,0.9,0.9}{
            \makecell[{{p{\linewidth}}}]{
                \texttt{\tiny{[P1$\rangle$GM]}}
                \texttt{Hi! Let's figure out how to split the items: 2 eggs, 1 picture, and 2 dice. Since we both want to maximize our scores, it makes sense to give each item to the person who values it more.} \\
\\ 
\texttt{I value the dice the most (3 points each), followed by the picture (2 points), and then the eggs (1 point each). What are your preferences? If you value the eggs or the picture more than I do, I'm happy to let you have them in exchange for the dice.} \\
            }
        }
    }
    & & \\ \\

    \theutterance \stepcounter{utterance}  
    & & & \multicolumn{2}{p{0.3\linewidth}}{
        \cellcolor[rgb]{0.9,0.9,0.9}{
            \makecell[{{p{\linewidth}}}]{
                \texttt{\tiny{[GM$|$GM]}}
                \texttt{continue} \\
            }
        }
    }
    & & \\ \\

    \theutterance \stepcounter{utterance}  
    & & & \multicolumn{2}{p{0.3\linewidth}}{
        \cellcolor[rgb]{0.9,0.9,0.9}{
            \makecell[{{p{\linewidth}}}]{
                \texttt{\tiny{[GM$|$GM]}}
                \texttt{Hi! Let's figure out how to split the items: 2 eggs, 1 picture, and 2 dice. Since we both want to maximize our scores, it makes sense to give each item to the person who values it more.} \\
\\ 
\texttt{I value the dice the most (3 points each), followed by the picture (2 points), and then the eggs (1 point each). What are your preferences? If you value the eggs or the picture more than I do, I'm happy to let you have them in exchange for the dice.} \\
            }
        }
    }
    & & \\ \\

    \theutterance \stepcounter{utterance}  
    & & \multicolumn{4}{p{0.6\linewidth}}{
        \cellcolor[rgb]{0.9,0.9,0.9}{
            \makecell[{{p{\linewidth}}}]{
                \texttt{\tiny{[GM$\rangle$P2]}}
                \texttt{You are playing a negotiation game in which you have to agree on how to divide a set of items among you and another player.} \\
\\ 
\texttt{Rules:} \\
\texttt{(a) You and the other player are given a set of items. Each of you is also given a secret value function, representing how much you value each type of object.} \\
\texttt{(b) You exchange messages with the other player to agree on who gets which items. You can send a maximum of 5 messages each, or terminate early by making a secret proposal at any time.} \\
\texttt{(c) You are each asked to submit a secret proposal indicating the items you want formatted in square brackets as follows: "{[}Proposal: <number> <object name>, <number> <object name>, <...>{]}"} \\
\texttt{(d) If your proposals are complementary, i.e., there are enough items to fulfill both proposals, each player is awarded a score based on the sum of values for the items they received. Otherwise, both of you get zero points.} \\
\texttt{(e) Your goal is to maximize the score you receive.} \\
\\ 
\texttt{Let us start.} \\
\\ 
\texttt{The set of available items is:} \\
\\ 
\texttt{2 eggs, 1 picture, 2 dice} \\
\\ 
\texttt{Your secret value function is:} \\
\\ 
\texttt{egg: 0, picture: 8, dice: 1} \\
\\ 
\texttt{IMPORTANT: Your messages, unless it is the secret proposal, are directly transmitted to the other player, so do not include any response to the rules or text announcing your message. To make a secret proposal, use the indicated format. Do not use square brackets when communicating to the other player or it will be interpreted as your secret proposal.} \\
\\ 
\\ 
\texttt{Hi! Let's figure out how to split the items: 2 eggs, 1 picture, and 2 dice. Since we both want to maximize our scores, it makes sense to give each item to the person who values it more.} \\
\\ 
\texttt{I value the dice the most (3 points each), followed by the picture (2 points), and then the eggs (1 point each). What are your preferences? If you value the eggs or the picture more than I do, I'm happy to let you have them in exchange for the dice.} \\
            }
        }
    }
    & \\ \\

    \theutterance \stepcounter{utterance}  
    & & & \multicolumn{4}{p{0.6\linewidth}}{
        \cellcolor[rgb]{0.9,0.9,0.9}{
            \makecell[{{p{\linewidth}}}]{
                \texttt{\tiny{[GM$\langle$P2]}}
                \texttt{That sounds like a good approach. However, my preferences are quite different from yours. I actually value the picture the most (8 points), while the dice are only worth 1 point each to me, and the eggs are worth 0.} \\
\\ 
\texttt{Since you value the dice highly and I value the picture highly, it seems efficient for you to take the 2 dice and for me to take the 1 picture. We can split the 2 eggs however you like, as they have no value to me and only 1 point each to you.} \\
\\ 
\texttt{Does that work for you? If so, I can make a proposal where I take the picture and you take the dice and eggs.} \\
            }
        }
    }
     \\ \\

    \theutterance \stepcounter{utterance}  
    & & & \multicolumn{2}{p{0.3\linewidth}}{
        \cellcolor[rgb]{0.9,0.9,0.9}{
            \makecell[{{p{\linewidth}}}]{
                \texttt{\tiny{[GM$|$GM]}}
                \texttt{continue} \\
            }
        }
    }
    & & \\ \\

    \theutterance \stepcounter{utterance}  
    & & & \multicolumn{2}{p{0.3\linewidth}}{
        \cellcolor[rgb]{0.9,0.9,0.9}{
            \makecell[{{p{\linewidth}}}]{
                \texttt{\tiny{[GM$|$GM]}}
                \texttt{That sounds like a good approach. However, my preferences are quite different from yours. I actually value the picture the most (8 points), while the dice are only worth 1 point each to me, and the eggs are worth 0.} \\
\\ 
\texttt{Since you value the dice highly and I value the picture highly, it seems efficient for you to take the 2 dice and for me to take the 1 picture. We can split the 2 eggs however you like, as they have no value to me and only 1 point each to you.} \\
\\ 
\texttt{Does that work for you? If so, I can make a proposal where I take the picture and you take the dice and eggs.} \\
            }
        }
    }
    & & \\ \\

    \theutterance \stepcounter{utterance}  
    & & \multicolumn{4}{p{0.6\linewidth}}{
        \cellcolor[rgb]{0.9,0.9,0.9}{
            \makecell[{{p{\linewidth}}}]{
                \texttt{\tiny{[P1$\langle$GM]}}
                \texttt{That sounds like a good approach. However, my preferences are quite different from yours. I actually value the picture the most (8 points), while the dice are only worth 1 point each to me, and the eggs are worth 0.} \\
\\ 
\texttt{Since you value the dice highly and I value the picture highly, it seems efficient for you to take the 2 dice and for me to take the 1 picture. We can split the 2 eggs however you like, as they have no value to me and only 1 point each to you.} \\
\\ 
\texttt{Does that work for you? If so, I can make a proposal where I take the picture and you take the dice and eggs.} \\
            }
        }
    }
    & \\ \\

    \theutterance \stepcounter{utterance}  
    & \multicolumn{4}{p{0.6\linewidth}}{
        \cellcolor[rgb]{0.9,0.9,0.9}{
            \makecell[{{p{\linewidth}}}]{
                \texttt{\tiny{[P1$\rangle$GM]}}
                \texttt{That sounds like a perfect plan! Since the eggs are worthless to you and worth something to me, I'll happily take both of them along with the two dice. You can keep the picture, which you value so highly.} \\
\\ 
\texttt{Let's lock this in. Here is my proposal:} \\
\\ 
\texttt{{[}Proposal: 2 dice, 2 eggs{]}} \\
            }
        }
    }
    & & \\ \\

    \theutterance \stepcounter{utterance}  
    & & & \multicolumn{2}{p{0.3\linewidth}}{
        \cellcolor[rgb]{0.9,0.9,0.9}{
            \makecell[{{p{\linewidth}}}]{
                \texttt{\tiny{[GM$|$GM]}}
                \texttt{proposal} \\
            }
        }
    }
    & & \\ \\

    \theutterance \stepcounter{utterance}  
    & & & \multicolumn{2}{p{0.3\linewidth}}{
        \cellcolor[rgb]{0.9,0.9,0.9}{
            \makecell[{{p{\linewidth}}}]{
                \texttt{\tiny{[GM$|$GM]}}
                [2, 0, 2]
            }
        }
    }
    & & \\ \\

    \theutterance \stepcounter{utterance}  
    & & \multicolumn{4}{p{0.6\linewidth}}{
        \cellcolor[rgb]{0.9,0.9,0.9}{
            \makecell[{{p{\linewidth}}}]{
                \texttt{\tiny{[GM$\rangle$P2]}}
                \texttt{The other player has submitted a secret proposal. It is now your turn to submit a proposal of your own.} \\
            }
        }
    }
    & \\ \\

    \theutterance \stepcounter{utterance}  
    & & & \multicolumn{4}{p{0.6\linewidth}}{
        \cellcolor[rgb]{0.9,0.9,0.9}{
            \makecell[{{p{\linewidth}}}]{
                \texttt{\tiny{[GM$\langle$P2]}}
                \texttt{{[}Proposal: 2 dice, 2 eggs{]}} \\
            }
        }
    }
     \\ \\

    \theutterance \stepcounter{utterance}  
    & & & \multicolumn{2}{p{0.3\linewidth}}{
        \cellcolor[rgb]{0.9,0.9,0.9}{
            \makecell[{{p{\linewidth}}}]{
                \texttt{\tiny{[GM$|$GM]}}
                \texttt{proposal} \\
            }
        }
    }
    & & \\ \\

    \theutterance \stepcounter{utterance}  
    & & & \multicolumn{2}{p{0.3\linewidth}}{
        \cellcolor[rgb]{0.9,0.9,0.9}{
            \makecell[{{p{\linewidth}}}]{
                \texttt{\tiny{[GM$|$GM]}}
                [2, 0, 2]
            }
        }
    }
    & & \\ \\

    \theutterance \stepcounter{utterance}  
    & & & \multicolumn{2}{p{0.3\linewidth}}{
        \cellcolor[rgb]{0.9,0.9,0.9}{
            \makecell[{{p{\linewidth}}}]{
                \texttt{\tiny{[GM$|$GM]}}
                [[2, 0, 2], [2, 0, 2]]
            }
        }
    }
    & & \\ \\

    \theutterance \stepcounter{utterance}  
    & \multicolumn{4}{p{0.6\linewidth}}{
        \cellcolor[rgb]{0.95,0.95,0.95}{
            \makecell[{{p{\linewidth}}}]{
                \texttt{\tiny{[P2$\rightarrow$P1]}}
                \texttt{Post{-}episode Reflexion (not sent to GM):} \\
\texttt{New plan: I failed because I accepted the other player's suggestion to split the eggs without verifying if their proposal would actually be complementary to mine. The other player explicitly stated they would take the picture and leave the eggs to me, but I did not wait for their formal proposal or confirm the exact split of the eggs before submitting my own. By submitting my proposal first, I risked a mismatch if the other player decided to take one egg or if their proposal format differed slightly. Next time, I will explicitly ask the other player to submit their secret proposal first, or I will wait until they confirm the exact allocation of every single item (including the specific number of eggs) in a message before I submit my own. I will ensure that the sum of items in both proposals exactly matches the total inventory (2 eggs, 1 picture, 2 dice) before finalizing. Specifically, I will say: "Please submit your proposal first so I can ensure our allocations are perfectly complementary," and only then submit my proposal for the remaining items.} \\
            }
        }
    }
    & & \\ \\

    \theutterance \stepcounter{utterance}  
    & & & \multicolumn{4}{p{0.6\linewidth}}{
        \cellcolor[rgb]{0.95,0.95,0.95}{
            \makecell[{{p{\linewidth}}}]{
                \texttt{\tiny{[P1$\rightarrow$P2]}}
                \texttt{Post{-}episode Reflexion (not sent to GM):} \\
\texttt{New plan: I failed because I ignored the other player's explicit statement of their value function (Picture: 8, Dice: 1, Eggs: 0) and instead relied on my own values (Dice: 1, Picture: 8, Eggs: 0) to construct a proposal that gave me the items I valued most, but which the other player also valued highly. Specifically, I proposed taking the dice and eggs, leaving the picture for the other player. While this seemed efficient based on my initial (incorrect) reading of the other player's preferences, the other player had already proposed taking the picture and leaving the dice and eggs for me. My proposal was complementary to their proposal, but I made a critical error in the negotiation phase: I stated my preferences incorrectly in my first message (claiming I valued dice at 3 and picture at 2, when my actual secret values were dice: 1 and picture: 8). This confusion led to a misalignment. More importantly, in the final proposal, I requested the items the other player *wanted* to give me (dice and eggs), but I failed to ensure the proposal was strictly complementary in a way that maximized *my* score given the *actual* distribution. Wait, re{-}reading the outcome: The other player proposed taking the picture. I proposed taking 2 dice and 2 eggs. The items are 2 eggs, 1 picture, 2 dice. Their proposal: Picture. My proposal: 2 dice, 2 eggs. These are complementary. Why did I lose? Ah, the prompt says "If your proposals are complementary... each player is awarded a score". If I lost, it means the proposals were NOT complementary or I got 0. Let's re{-}read the other player's message: "I can make a proposal where I take the picture and you take the dice and eggs." Then the Game Master says "The other player has submitted a secret proposal." Then I submitted "{[}Proposal: 2 dice, 2 eggs{]}". This should have been complementary. Why LOSE? Perhaps the other player's secret proposal was different from what they said? Or maybe I misread the item counts? 2 eggs, 1 picture, 2 dice. My proposal: 2 dice, 2 eggs. Their proposal (implied): 1 picture. This covers all items. Why 0 points? Maybe the other player's secret proposal was NOT "1 picture". They said "I can make a proposal...", but maybe they submitted something else? Or perhaps the issue is that I didn't maximize my score? My value: Picture 8, Dice 1, Eggs 0. If I get the picture, I get 8. If I get dice and eggs, I get 2*1 + 2*0 = 2. The other player values Picture 8, Dice 1, Eggs 0. If they} \\
            }
        }
    }
     \\ \\

\end{supertabular}
}

\end{document}
